% ============================================================
%  Part I — Sounds and spelling
% ============================================================
\gpart{I}{Sounds \& Spelling}

% ------------------------------------------------------------
\gsec{The alphabet and stress}

\gintro{German spelling is close to phonetic: once you know the rules in this
part, a written word tells you how it sounds. The approximations used here are
based on British English.}

\tcap{The alphabet}

\begin{dtbl}[\footnotesize]{C{0.55}W{1.45}SC{0.55}W{1.45}SC{0.55}W{1.45}SC{0.55}W{1.45}}
\rowcolor{Ink} \hd{} & \hd{Name} & \hd{} & \hd{Name} & \hd{} & \hd{Name} & \hd{} & \hd{Name} \\ \hdrule
A & ah   & H & hah  & O & oh   & V & fau \\
B & beh  & I & ih   & P & peh  & W & weh \\
C & tseh & J & jott & Q & kuh  & X & iks \\
D & deh  & K & kah  & R & err  & Y & üpsilon \\
E & eh   & L & ell  & S & ess  & Z & tsett \\
F & eff  & M & emm  & T & teh  & Ä & ah-Umlaut \\
G & geh  & N & enn  & U & uh   & Ö & oh-Umlaut \\
  &      &   &      & Ü & uh-Umlaut & ß & Eszett \\
\end{dtbl}

\begin{gnote}
\textbf{ß} is called \textit{Eszett} or \textit{scharfes S} and always sounds
like a sharp \emph{s}. It never begins a word and is replaced by \textit{ss} in
Switzerland. The umlauts are alphabetised as though they were plain
\textit{a, o, u}.
\end{gnote}

\tcap{Where the stress falls}

\begin{dtbl}[\footnotesize]{W{0.7}W{1.3}}
\rowcolor{Ink} \hd{Rule} & \hd{Example} \\ \hdrule
Most native words  & on the first syllable: \textbf{Ar}beit, \textbf{Kin}der \\
Separable prefix   & on the prefix: \textbf{auf}stehen, \textbf{an}kommen \\
Inseparable prefix & on the stem: ver\textbf{ste}hen, be\textbf{zah}len \\
Verbs in -ieren    & on the \textit{-ier}: stu\textbf{die}ren \\
Many loanwords     & near the end: Universi\textbf{tät}, Na\textbf{tion} \\
\end{dtbl}

\begin{gnote}
Stress alone distinguishes \textbf{über}setzen \emph{to ferry across} from
über\textbf{set}zen \emph{to translate}, which is why the dual-prefix verbs on
page 26 split into two groups.
\end{gnote}

% ------------------------------------------------------------
\gsec{Vowels}

\begin{dtbl}[\footnotesize]{W{0.55}W{1.15}W{1.3}}
\rowcolor{Ink} \hd{Letter} & \hd{Sounds like} & \hd{Example} \\ \hdrule
a  & short \emph{a} in \emph{cat}, long \emph{a} in \emph{father} & Mann, Vater \\
e  & \emph{e} in \emph{bed}, or \emph{ay} in \emph{say}   & Bett, gehen \\
i  & \emph{i} in \emph{sit}, or \emph{ee} in \emph{see}   & bitte, Tiger \\
o  & \emph{o} in \emph{pot}, or \emph{o} in \emph{go}     & Sonne, Brot \\
u  & \emph{u} in \emph{put}, or \emph{oo} in \emph{boot}  & Mutter, Buch \\
ä  & \emph{e} in \emph{bed}                               & Männer \\
ö  & \emph{ur} in \emph{fur}, lips rounded                & schön \\
ü  & \emph{ee} said with rounded lips                     & über \\
\end{dtbl}

\tcap{Vowel combinations}

\begin{dtbl}[\footnotesize]{W{0.55}W{1.15}W{1.3}}
\rowcolor{Ink} \hd{Letters} & \hd{Sounds like} & \hd{Example} \\ \hdrule
ei, ai & \emph{eye}              & mein, Mai \\
ie     & \emph{ee} in \emph{see} & Liebe \\
eu, äu & \emph{oy} in \emph{boy} & neu, Häuser \\
au     & \emph{ow} in \emph{how} & Haus \\
\end{dtbl}

\begin{gnote}
\textit{ei} and \textit{ie} are the pair most often confused. Each says the name
of its \textbf{second} letter: \textit{ei} sounds like \emph{I}, \textit{ie} like
\emph{E}.
\end{gnote}

\tcap{Long or short?}

\begin{dtbl}[\footnotesize]{W{0.7}W{1.3}}
\rowcolor{Ink} \hd{The vowel is} & \hd{When} \\ \hdrule
long  & before \textit{h}: \textit{Zahn}, \textit{Uhr}, \textit{gehen} \\
long  & when doubled: \textit{Boot}, \textit{Tee}, \textit{Saal} \\
long  & before a single consonant: \textit{Tag}, \textit{Brot} \\
short & before a double consonant: \textit{kommen}, \textit{Mutter} \\
short & before most consonant clusters: \textit{Kind}, \textit{Herbst} \\
\end{dtbl}

% ------------------------------------------------------------
\gsec{Consonants}

\begin{dtbl}[\footnotesize]{W{0.5}W{1.2}W{1.3}}
\rowcolor{Ink} \hd{Letters} & \hd{Sounds like} & \hd{Example} \\ \hdrule
w      & English \emph{v}                        & Wasser, wo \\
v      & English \emph{f}                        & Vater, viel \\
v      & English \emph{v}, in loanwords          & Vase, Klavier \\
z      & \emph{ts} in \emph{cats}                & Zeit, zwei \\
j      & English \emph{y}                        & ja, Jahr \\
s      & English \emph{z} before a vowel         & Sonne, sagen \\
ss, ß  & sharp \emph{s}                          & Wasser, Fuß \\
sch    & \emph{sh}                               & Schule \\
sp, st & \emph{shp}, \emph{sht} at the start of a word & sprechen, Stadt \\
ch     & throaty, after a, o, u, au              & Bach, Buch \\
ch     & soft, like \emph{h} in \emph{huge}, elsewhere & ich, Milch \\
chs    & \emph{ks}                               & Fuchs, sechs \\
ng     & \emph{ng} in \emph{singer}, no hard g   & singen \\
pf     & both letters sounded                    & Pfanne \\
qu     & \emph{kv}                               & Quelle \\
r      & from the back of the throat             & rot \\
-er    & like \emph{uh} at the end of a word     & Vater, Mutter \\
b, d, g & harden to \emph{p, t, k} at the end of a word & Hund, Tag, gelb \\
h      & silent after a vowel, lengthening it    & gehen, Uhr \\
th, ph & \emph{t}, \emph{f}                      & Theater, Physik \\
\end{dtbl}

\begin{gnote}
The two values of \textit{ch} are decided by the vowel in front of it. After
\textit{a, o, u, au} it is the rasping sound in \textit{Bach}; after anything
else, including all three umlauts, it is the soft sound in \textit{ich}.
\end{gnote}

\begin{gnote}
Final \textit{b}, \textit{d} and \textit{g} harden, so \textit{Hund} sounds like
\emph{hunt} and \textit{Tag} like \emph{tak}. The soft sound returns as soon as
an ending follows: \textit{Hunde}, \textit{Tage}.
\end{gnote}
